\documentclass[12pt, a4paper]{article}

% --- Pacotes Básicos ---
\usepackage[utf8]{inputenc}
\usepackage[T1]{fontenc}
\usepackage[main=brazilian, english]{babel}
\usepackage{indentfirst}
\usepackage{geometry}
\geometry{left=3cm, top=3cm, right=2cm, bottom=2cm}

% --- Informações do Documento ---
\title{Diário De Bordo: Aula 3\\Iniciação à Prática de Extensão}
\author{Nome: Erickson Giesel Müller
\\ Fase: 6° Semestre Ciência da Computação}
\date{\today}

\begin{document}
\maketitle
\thispagestyle{empty} % primeira página não tem número.

\section*{Projetos de extensão na linha de tecnologias educacionais}
O construtivismo de Jean Piaget inspiriou Seymour Papert a formar o construcionismo por volta de 1950, essa é a base da teconologia educacional. Quando Papert visitou a Unicamp, em 1970, conheceu José Armando Valente, levou-o ao MIT onde concluiu seu mestrado e doutorado e criou o Media Lab com foco em tecnologias educacionais.

O AppInventor é a plataforma que usaremos para criar o projeto de aplicação mobile da CCR, e vamos seguir uma das 3 teorias de conhecimento: construcionismo, conectivismo e teoria da carga cognitiva. Em 2020, as tecnologias educacionais tiveram destaque por conta da pandemia.

\section*{Construcionismo}
A teoria contrucionista parte do princípio que as pessoas aprendem melhor quando estão envolvidas de maneira ativa na cosntrução de coisas no mundo. A teoria incentiva o aluno a identificar o próprio erro e que, ao consertar esse erro, o aprendizado é engessado (método \textit{debug}).

\section*{Conectivismo}
Em 2003, George Siemens criou a teoria que nós estamos em rede e só conseguimos avançar socialmente quando estamos em conexão. O conhecimento vai se apoiar em diversidade de opiniões, as pessoas só aprendem com pessoas que pensam diferente delas. Aprendizagem mais horizontal. Aprendizagem com dispositivos não humanos.

\section*{Teoria da Carga Cognitiva}
A carga cognitiva está relacionada com a quantidade de informação que uma pessoa pode armazenar na memória. O nosso organismo é sensível ao recebimento de informações, o conhecimento prévio é armazenado na memória de longo prazo.

As cargas cognitivas são: carga intrínseca e carga estranha. A carga intrínseca se refere às complexidades inerentes ao conteúdo em estudo e a carga estranha se remete ao modo como o conteúdo em estudo é apresentado ao aprendiz aos excessos e supérfluos adotados pelo instrutor. Esta última é mais moldável para o aprendizado que será apresentado ao aluno pelo professor.

\section*{Informática na Educação no Brasil}
	\begin{itemize}
		\item Anos 70: Valente(Unicamp), Fagundes(UFRGS).
		\item Anos 80: Seminários em Brasília e Bahia, Projeto EDUCOM, com o fim do governo militar a educação mudou a forma como era vista pela sociedade. PRONINFE.
		\item Anos 90: Internet no Brasil, PROINFO, Telecursos, TVEscola.
		\item Anos 2000: EAD no ensino superior, objetos educacionais, repositórios como RIVED, UAB.
		\item Anos 2010: PROUCA, P\&D
	\end{itemize}

\section*{Fundo de Universalização dos Serviços de Telecomunicações}
	O FUST dá apoio a serviços de teleeducação e telemedicina. O objetivo de destinar recursos para cobrir os custos do cumprimento das obrigações de universalização dos serviços de telecomunicações que não podem ser recuperados através da exploração eficiente do serviço.


\end{document}
